\section{Deterministic Smart Contracts}
\label{sec:vm}
We treat programmability as a \emph{supporting capability}, not a scientific contribution of this
paper: the core results concern fragmented history and the storage-coupled consensus, and a reader
interested only in those may skip this section. We include it because a storage-scarce ledger is of
little use if it cannot run application logic, and because the contract layer must obey the same
commitment discipline (signed roots, deterministic execution) as the rest of the system; we
therefore keep the description brief.

Contracts execute on a deterministic stack virtual machine with metered gas. Determinism is a hard
requirement: every node that runs the same bytecode on the same input and storage must reach the
same result, or consensus over contract state breaks. The VM therefore excludes every source of
nondeterminism, namely wall-clock time, floating point, and randomness, and bounds execution by a
gas budget and a step limit so a contract cannot loop forever or exhaust memory.

\textbf{Machine model.} The VM is a stack machine over $64$-bit words with a per-contract
key/value store (uint64 to uint64). Opcodes fall into the classes of Table~\ref{tab:opcodes}.
Storage writes are the expensive operation, reflecting their real cost (they enlarge the state
every node keeps), so gas pricing discourages gratuitous writes. Cross-contract calls are allowed
but bounded by a call-depth limit, which caps reentrancy and recursion. Because a $64$-bit word
cannot hold a $20$-byte address or a string, a byte-storage layer holds variable-length blobs
(addresses, token identifiers, URIs) sourced from calldata, with opcodes to store, compare, hash,
and return them; this is how a contract can manipulate full addresses despite the uint64 core.

\begin{table}[t]
\centering
\caption{VM opcode classes (representative, not exhaustive).}
\label{tab:opcodes}
\small
\begin{tabular}{ll}
\toprule
class & opcodes / role \\
\midrule
arithmetic    & add/sub/mul (overflow-checked), div/mod (revert on $0$) \\
logic/compare & and, or, not, eq, lt, gt, iszero \\
stack         & push, pop, dup, swap \\
control flow  & jump, jumpi, stop, return, revert \\
storage       & sload, sstore, mix (mapping-slot hash) \\
environment   & caller, value, self, calldata, calldata-len \\
byte layer    & bstore, beq, bhash, caller-bytes, return-bytes \\
events / calls& log (committed in LogsRoot), call (depth-bounded) \\
\bottomrule
\end{tabular}
\end{table}

\textbf{The SHL language.} Writing raw bytecode is error-prone, so contracts are written in SHL, a
small Solidity-like language compiled through a lexer, parser, and code generator. SHL offers named
\texttt{state} variables and mappings (each assigned a fixed storage slot by the compiler, so there
are no hand-managed magic offsets), functions with selector dispatch, \texttt{require} and
\texttt{revert} with full state rollback, and overflow-checked arithmetic by default. A function
selector is derived from the function name and placed in the call's first calldata word; the
compiler emits a dispatcher that routes to the matching function and reverts on an unknown selector.
Listing~\ref{lst:shl} shows a fungible-token transfer: the \texttt{require} guard reverts the entire
call (restoring all storage) if the sender is short, and both updates are overflow-checked, so a
short and readable contract is also a safe one.

\begin{lstlisting}[caption={A safe ERC-20-style transfer in SHL.},label={lst:shl}]
state balances;          // mapping holder -> amount
state totalSupply;

fn transfer(to, amount) {
  require(balances[msg.sender] >= amount);
  balances[msg.sender] = balances[msg.sender] - amount; // checked
  balances[to] = balances[to] + amount;                 // checked
  emit(TRANSFER, msg.sender, to, amount);
}
fn balanceOf(who) { return balances[who]; }
\end{lstlisting}

\textbf{Scope.} SHL is deliberately a faithful \emph{subset} of the Solidity model rather than a
clone. It omits $256$-bit integers, contract inheritance, and the ABI, trading generality for a VM
small enough to audit and reason about, which suits a system that has not yet undergone external
security review (\S\ref{sec:disc}). Events emitted via \texttt{log} are committed through the block
header's logs root, so a verifier can later prove that an event occurred without trusting the node
that serves it, the same commitment discipline used for shards and transactions.
